% 天蝎座（综述）
% license CCBYSA3
% type Wiki

本文根据 CC-BY-SA 协议转载翻译自维基百科\href{https://en.wikipedia.org/wiki/Scorpius}{相关文章}。

天蝎座（Scorpius）是一个黄道星座，位于南天半球，处在银河系中心附近，其西侧为天秤座，东侧为人马座。天蝎座是一组极为古老的星座，其被认识的历史早于希腊文化$^{[2]}$；它也是公元$^{[2]}$世纪希腊天文学家托勒密所确认的 48 个星座之一。
\subsection{显著特征}  
\subsubsection{恒星}
\begin{figure}[ht]
\centering
\includegraphics[width=8cm]{./figures/2cc53693047d69c0.png}
\caption{肉眼所见的天蝎座星座（已绘制星座连线）。} \label{fig_Scorpi_1}
\end{figure}
天蝎座包含许多明亮的恒星，其中包括心宿二 Antares（$\alpha$ Sco），意为“火星的对手”，因其显著的红色光泽而得名；$\beta_1$ Sco（Graffias 或 Acrab），一套三星系统；$\delta$ Sco（Dschubba，“额头”）；$\theta$ Sco（Sargas，名称源自苏美尔文化 $^{[3]}$）；$\nu$ Sco（Jabbah）；$\xi$ Sco；$\pi$ Sco（Fang）；$\sigma$ Sco（Alniyat）；以及 $\tau$ Sco（Paikauhale）。

位于蝎子弯曲尾部末端的是$\lambda$Sco（Shaula）与$\upsilon$Sco（Lesath），两者的名称均意为“毒刺”。由于它们彼此距离很近，$\lambda$Sco与$\upsilon$Sco 有时被合称为“猫眼” $^{[4]}$。

天蝎座中这些明亮恒星组成的形状类似于码头工人所用的钩子。其中大多数恒星都是最近的 OB 星协——天蝎—半人马星协（Scorpius--Centaurus）的巨大质量成员$^{[5]}$。

恒星$\delta$Sco 在长期保持稳定的 2.3 等亮度之后，于 2000 年 7 月在数周内突然增亮至 1.9 等。此后它成为一颗变星，亮度在 2.0 等与 1.6 等之间波动$^{[6]}$。这意味着在最亮时，它是天蝎座中第二亮的恒星。
\begin{figure}[ht]
\centering
\includegraphics[width=8cm]{./figures/a6e325eb5a989396.png}
\caption{按距离（红--绿三维视图）显示的天蝎座恒星分布，以及各恒星亮度（以恒星大小表示）} \label{fig_Scorpi_2}
\end{figure}
天蝎座 U 星（U~Scorpii）是目前已知爆发周期最短的新星，其周期约为 10 年$^{[7]}$。

天蝎座 AH 星是一颗红超巨星，也是已知体积最大的恒星之一，其半径约为太阳的 1400 倍。同时它也是一颗高光度恒星，亮度约为太阳的 340000 倍$^{[8]}$，但由于其视星等在 6.5 至 9.6 之间变化$^{[9]}$，肉眼通常无法直接观测到。

由$\omega^1$天蝎座 与$\omega^2$天蝎座 组成的一对近邻恒星是一对视双星，可用肉眼分辨。其中一颗是黄色巨星$^{[10]}$，另一颗则是天蝎—半人马星协中的一颗 B 型蓝色恒星$^{[11]}$。

曾被指定为$\gamma$天蝎座的恒星（尽管其实际位于天秤座边界之内）如今被称为$\sigma$天秤座。此外，在古希腊时期，整个天秤座曾被视为天蝎座的钳爪（Chelae~Scorpionis）；而在稍后的希腊时代，这些最西侧的恒星被重新诠释为由正义女神 Astraea（由相邻的处女座所代表）高举的一副天平。天秤座作为独立星座的划分在古希腊或古罗马时期被正式确立$^{[12]}$。
\subsubsection{深空天体}
\begin{figure}[ht]
\centering
\includegraphics[width=6cm]{./figures/36a84b2fd26c01cc.png}
\caption{天蝎座与银河系，同画面中可见心宿二（Antares）附近的 M4 与 M80，画面中央下方的 M6 与 M7，画面顶部的 NGC\,6124，以及画面中央正上方的 NGC\,6334。} \label{fig_Scorpi_3}
\end{figure}
由于其位置横跨银河系，这个星座包含了大量深空天体，例如疏散星团梅西耶~6（蝴蝶星团）与梅西耶~7（托勒密星团），位于 $\zeta^2$ 天蝎座附近的 NGC~6231，以及球状星团梅西耶~4 与梅西耶~80。

梅西耶~80（NGC~6093）是一座视星等为 7.3 的球状星团，距离地球约 33000 光年。它属于 Shapley~II 级球状星团；这一分类表明其核心高度集中且密度极高。M80 由查尔斯·梅西耶于$^{[1781]}$年发现。1860 年，阿图尔·冯·奥沃斯在此发现了新星天蝎座~T 星，这是一次相当罕见的发现$^{[13]}$。

NGC~6302，又称“甲虫星云”，是一座双极行星状星云。NGC~6334，又称“猫爪星云”，是一片发射星云及恒星形成区域。
\begin{figure}[ht]
\centering
\includegraphics[width=6cm]{./figures/5de286b7a5f3002e.png}
\caption{天蝎座的核心区域。画面中央偏左可见 M4。$\rho$蛇夫座分子云复合体的部分区域被心宿二（Antares）及其周围的恒星所照亮。} \label{fig_Scorpi_4}
\end{figure}
\subsection{神话}
\begin{figure}[ht]
\centering
\includegraphics[width=6cm]{./figures/adbb7e7cf91cc44d.png}
\caption{天蝎座在《乌拉尼娅之镜》中所描绘的形象；《乌拉尼娅之镜》是一套约于 1825 年在伦敦出版的星座卡片。} \label{fig_Scorpi_5}
\end{figure}
在希腊神话中，与天蝎座相关的若干神话都将其归因于猎户座。根据其中一种说法，猎户座向女神阿尔忒弥斯及其母亲勒托夸口，说他将杀死地球上的所有动物。阿尔忒弥斯和勒托于是派出一只蝎子去杀死猎户座$^{[14]}$。他们的战斗引起了宙斯的注意，宙斯将双方都升上天空，以此提醒凡人要抑制过度的傲慢。在另一种版本的神话中，是阿尔忒弥斯的孪生兄弟阿波罗派出蝎子杀死猎户座，因为猎人因承认女神比自己更强而赢得了她的青睐。宙斯将猎户座和蝎子升上天空之后，前者每到冬季便在天空中追逐，而每到夏季，当蝎子出现时便逃离。在这两种版本中，都是阿尔忒弥斯请求宙斯将猎户座升上天空。

在另一则不涉及猎户座的希腊神话中，天上的蝎子曾在法厄同驾驶其父赫利俄斯的太阳战车时与其相遇$^{[15]}$。
\subsubsection{起源}  
巴比伦人称这一星座为 MUL.GIR.TAB——“蝎子”；该名称可以直译为“具有燃烧之刺的生物”$^{[16]}$。

在一些古老的记述中，天秤座被视为天蝎的钳爪。天秤座在巴比伦文献中被称为蝎之钳（zibānītu，参见阿拉伯语 zubānā），在希腊文中则称为 χηλαι$^{[17]}$。
\subsubsection{占星学}  
主条目：Scorpio（占星学）  
西方占星学中的天蝎座与天文学上的天蝎座星座并不相同。从天文学角度看，太阳仅在 11 月 23 日至 11 月 28 日这 6 天内位于国际天文学联合会定义的天蝎座边界之内。差异的很大一部分源于蛇夫座的存在，而该星座仅被少数占星家采用。天蝎座在印度占星体系中对应的宿为 Anuradha、Jyeshtha 和 Mula$^{[citation\ needed]}$。
\subsection{文化}
\begin{itemize}
\item 印度尼西亚的爪哇人将这一星座称为 Banyakangrem（意为“孵卵的天鹅”）或 Kalapa Doyong（意为“倾斜的椰子树”），这是因为其形状上的相似性。
\item 在夏威夷，天蝎座被称为半神毛伊（Māui）的鱼钩，称作 Maui’s Fishhook，或 Ka Makau Nui o Māui（意为“毛伊的大鱼钩”），而这把鱼钩的名字是 Manaiakalani。
\item 天蝎座曾被分为两个星群，并被布吉斯（Bugis）水手用于航海导航。天蝎座的北部（α、β、γ 或 σ 天秤座、δ、ε、ζ、μ、σ 与 τ 天蝎座）被称为 bintoéng lambarué，意为“鳐鱼之星”；天蝎座的南部（η、θ、ι、κ、λ 与 ν 天蝎座）被称为 bintoéng balé mangngiwéng，意为“鲨鱼之星”。
\item 在巴西国旗上，天蝎座代表巴西的东北地区，但巴伊亚州除外，其归属于南十字座。
\end{itemize}
\subsection{另见}
\begin{itemize}
\item 天蝎座（中国天文学）
\end{itemize}
\subsection{参考文献}
\begin{enumerate}
\item 柯克帕特里克，J·戴维；马罗科，费德里科等（2024年4月）。“基于全天空20秒差距范围内约3600颗恒星和棕矮星普查的初始质量函数”. 天体物理学杂志增刊系列. 271 (2): 55. arXiv:2312.03639. 文献标识码:2024ApJS..271...55K. doi:10.3847/1538-4365/ad24e2.
\item 奈特，J.D.“天蝎座——海天之间的星座”. www.seasky.org。检索日期：2017-02-11.
\item 伯纳姆，罗伯特（1978-01-01）。伯纳姆天体手册：面向观测者的太阳系外宇宙指南。库瑞尔公司。ISBN 978-0-486-23673-5.
\item 弗雷德·沙夫（麦克米伦，1988年）《认识星空的40个夜晚：逐夜观星入门指南》，第79页（ISBN 9780805046687).
\item 普雷比施，T.；马马耶克，E.（2009）。“最近的OB协会：天蝎座—半人马座（天蝎座OB2）”。恒星形成区域手册. 2: 0. arXiv:0809.0407. 文献标识码:2008hsf2.book..235P.
\item “天蝎座δ星依旧风采依旧”。存档于2007年6月6日的原始页面。检索日期：2008年6月28日。
\item AAVSO：本季变星：天蝎座U星.
\item 阿罗约-托雷斯，B.；维特科夫斯基，M.；马尔凯德，J. M.；豪斯希尔德特，P. H. （2013-06-01）。“红超巨星天蝎座AH、仙后座UY和人马座KW的大气结构与基本参数”. 天文学与天体物理学. 554：A76。arXiv:1305.6179. 文献标识码:2013A&A...554A..76A. DOI:10.1051/0004-6361/201220920. ISSN 0004-6361.
\item 基斯，L. L.；萨博，G. M.；贝丁，T. R.（2006-11-01）。“红超巨星的变异性：脉动、长次级周期与对流噪声”. 皇家天文学会月报. 372 (4): 1721–1734. arXiv:astro-ph/0608438. Bibcode:2006MNRAS.372.1721K. doi:10.1111/j.1365-2966.2006.10973.x. ISSN 0035-8711.
\item 侯克，N.；史密斯-摩尔，M.（1988）密歇根HD恒星二维光谱类型目录。赤纬范围：-26°.0至-12°.0, 第4卷,Bibcode:1988mcts.book.....H.
\item 吉林斯基，E.；等（2006年3月），“天蝎座-半人马座星协中B型恒星的径向速度测量”，天文学与天体物理学, 448 (3): 1001–1006, arXiv:astro-ph/0601643, Bibcode:2006A&A...448.1001J, doi:10.1051/0004-6361:20041614, S2CID 17818058.
\item 德克尔，埃莉（2013）。阐释现象：古代与中世纪的天体制图。牛津：牛津大学出版社。第61–62. ISBN 978-0-19-960969-7.
\item 利维 2005, 第166–167页。
\item 荷马注释关于伊利亚特18.486 引用费雷西德斯
\item 天蝎座——传说与神话存档于2008-07-20 互联网档案馆
\item 伍尔福克，乔安娜（2011）。天蝎座。兰汉姆：泰勒贸易出版社。第81页。ISBN 978-1589795600.
\item 巴比伦星象学加文·怀特著，索拉里亚出版社，2008年，第175页。
\item 达尔焦尼，N（1984）。“普拉纳塔芒萨：爪哇农业历法——其在发展农村生活中的生物气候与社会文化功能”。环境学家. 4（S7）：15–18. 文献标识码:1984ThEnv...4S..15D. DOI:10.1007/BF01907286. S2CID 189914684.
\item "Jejak Langkah Astronomi di Indonesia". 2 January 2011.
\item "Hawaiian Astronomical Society, Constellations: Scorpius - The Scorpion who Killed Orion".
\item “夏威夷星线与恒星名称——星线3：马奈阿卡拉尼”.
\item 凯利，大卫·H.；米隆，尤金·F.；阿维尼，A.F.（2011）。探索古代星空：古代与文化天文学概览。纽约，纽约：施普林格出版社。第344页。ISBN 978-1-4419-7623-9.
\end{enumerate}
\begin{itemize}
\item 利维，大卫·H.（2005）。深空天体. 普罗米修斯图书. ISBN 1-59102-361-0.
\item 伊恩·里德帕斯和威尔·蒂里昂（2007）。恒星与行星指南, 科林斯，伦敦。ISBN 978-0-00-725120-9。普林斯顿大学出版社，普林斯顿。ISBN 978-0-691-13556-4.
\end{itemize}
\subsection{外部链接}
\begin{itemize}
\item 天蝎座 (类别)
\item 深空摄影指南：天蝎座
\item 可点击的天蝎座
\item 星图——天蝎座
\item 瓦尔堡学院图像数据库（天蝎座的中世纪及早期现代图像）
\end{itemize}